% Section 4: Forecasting for RL -- model-based planning (Story 2, part 1)
\subsection{Forecasting for RL: model-based planning}
\label{sec:model_based_planning}

Section~\ref{sec:rl_for_forecasting} kept the forecaster as the artifact.
This section and the next invert the relation. The forecaster becomes a
component inside a reinforcement-learning agent, and the artifact is a
policy. The boundary established in \S\ref{sec:madeka} marks the entry.
Once the agent's actions feed back into the data-generating process, the
reduction to supervised learning fails and the problem is genuinely
reinforcement learning.

The object the agent now needs is a forecaster of its own environment.
Given a state $s$ and an action $a$, it must predict the next state and
reward, $\widehat P(s', r \mid s, a)$. An economist will recognise this
object. Strip the action argument and it is the state-space model that
underlies the Kalman filter and the ARIMA family. Restore the action and
it is the controlled state-space model, the autoregressive system with
exogenous inputs that control engineers and macroeconometricians have
fitted for decades.\footnote{The state-space forecasting tradition is
surveyed in \citet{DurbinKoopman2012}. The controlled generalisation, in
which an exogenous input enters the transition equation, is standard in
both the control literature and the macroeconometric vector-autoregression
literature.} The recurrent state-space model at the core of the Dreamer
architecture below is the neural generalisation of exactly that object.
Two ingredients separate a world model from its classical ancestor. The
action enters the transition, and the model is trained against a value
functional rather than a likelihood. The first ingredient is structural.
The second is the sequential form of the bridge of
Section~\ref{sec:bridge}, and it is the subject of \S\ref{sec:vaml_ve}.
The empirical architectures that put the two ingredients together are the
subject of \S\ref{sec:world_models}.

\subsubsection{Value-aware model learning}
\label{sec:vaml_ve}

The bridge of Section~\ref{sec:bridge} has a sequential form. Estimating
the transition kernel $\widehat P$ by maximum likelihood treats every
component of the next-state distribution as equally informative. The
downstream planner cares about only one functional of it, namely
$(\widehat P V)(s, a) = \mathbb E_{s' \sim \widehat P(\cdot \mid s, a)}
[V(s')]$ for value functions $V$ in some candidate class. A value-aware
loss aligns the model-learning objective with that functional, in direct
analogy to the single-period decision-focused losses of \S\ref{sec:dfl}.

\citet{FarahmandVAML2017} formalise the alignment as the value-aware
model-learning (VAML) loss
\begin{equation}
\mathcal L_{\mathrm{VAML}}(\widehat P)
= \sup_{V \in \mathcal V} \bigl\| (\widehat P V) - (P^\star V) \bigr\|_\mu,
\label{eq:vaml_loss}
\end{equation}
where $\mathcal V$ is the candidate value class, $P^\star$ is the true
kernel, and $\| \cdot \|_\mu$ is an $L^2$ norm under a state-action
distribution $\mu$. They prove a finite-sample upper bound on the
value-suboptimality of the policy planned in $\widehat P$ that decomposes
into a model-approximation term, an estimation term that shrinks at the
parametric rate, and a complexity term in
$\mathcal V$.\footnote{An alternative regularity-restricted form is the
Wasserstein-VAML equivalence of \citet{AsadiWasserstein2018}. When
$\mathcal V$ is the class of $1$-Lipschitz value functions,
$\mathcal L_{\mathrm{VAML}}$ coincides with the $1$-Wasserstein distance
between $\widehat P(\cdot \mid s, a)$ and $P^\star(\cdot \mid s, a)$ by
Kantorovich-Rubinstein duality, which gives an optimal-transport reading
of value-aware learning.} \citet{FarahmandIterVAML2018} embeds
\eqref{eq:vaml_loss} inside approximate value iteration. At iteration $k$,
the model is fit to minimise the value-prediction error against the
current iterate $V_k$ rather than against the full class $\mathcal V$,
which turns the optimisation into a tractable regression and yields an
LSPI-style performance bound on the converged policy.

The full generalisation is \citet{GrimmVALEQ2020}'s value-equivalence
principle. Two models $\widehat P_1, \widehat P_2$ are value-equivalent
with respect to a policy set $\Pi$ and a value-function set $\mathcal V$
when $(T_{\widehat P_1}^\pi V)(s) = (T_{\widehat P_2}^\pi V)(s)$ for every
$\pi \in \Pi$ and $V \in \mathcal V$, where $T_{\widehat P}^\pi$ is the
Bellman operator under model $\widehat P$ and policy $\pi$. The
equivalence class shrinks monotonically as $\Pi$ and $\mathcal V$ grow,
collapsing to the singleton $\{P^\star\}$ in the limit of all policies and
all value functions. The principle subsumes Value Iteration Networks, the
Predictron, Value Prediction Networks, TreeQN, and the MuZero family of
\S\ref{sec:world_models} as instances that train a learned model against
narrow $(\Pi, \mathcal V)$ rather than against observation
likelihood.\footnote{\citet{GrimmPVE2021} sharpen the result for the
MuZero loss specifically. Their proper value equivalence relaxes the
equivalence relation to the $k$-step Bellman operator and bounds the
value suboptimality of the planned policy in terms of the model's loss on
a value-policy pair drawn from the planner. The bound formalises why
MuZero can plan well with models that fail at observation prediction.}

\citet{AyoubVTR2020} convert the value-aware idea into a regret
guarantee. In their value-targeted regression (VTR) algorithm, the model
at episode $k$ is fit by
\begin{equation}
\widehat P_k \in \arg\min_{P \in \mathcal P}
\sum_{t < k} \Bigl(
  (P V_t)(s_t, a_t) - V_t(s_{t+1})
\Bigr)^2,
\label{eq:vtr_loss}
\end{equation}
where $V_t$ is the most recent value-function iterate and $\mathcal P$ is
a linear-mixture model class. They prove a regret bound of
$\widetilde O\bigl(d \sqrt{H^3 T}\bigr)$, where $d$ is the model-class
dimension, $H$ is the horizon, and $T$ is the number of episodes. The
result is the first regret guarantee in which model fitting is driven by
a value-prediction loss rather than by a likelihood, and the rate matches
a known $\Omega(\sqrt{d H T})$ lower bound up to $H$ and logarithmic
factors. The value targets in \eqref{eq:vtr_loss} are themselves
policy-dependent and evolve through the learning loop, which is what
makes value-aware model learning irreducibly reinforcement learning
rather than the supervised special case of \S\ref{sec:madeka}.

Value-aware losses are not without pitfalls. \citet{VoelckerCalib2025}
show that the MuZero loss and several of its variants are uncalibrated in
the surrogate-loss sense of Section~\ref{sec:bridge}. The population
minimiser of the MuZero loss can be a transition kernel whose induced
value function differs from $V^\star$ even when $V^\star$ lies in the
function class. They propose a calibration correction that couples the
model class with a stochastic latent architecture and recovers the right
population minimiser. The result transfers the calibration question that
occupied Section~\ref{sec:bridge} into the sequential setting and ties the
value-aware programme back to the proper scoring-rule machinery.

\subsubsection{Dyna and the world-model line}
\label{sec:world_models}

The idea of using a learned forecaster to generate experience for a
planner is old. \citet{sutton1990} introduced the Dyna architecture as
an integrated loop of three operations executed at every time step.
The agent acted in the real environment and collected a transition. It
updated a learned dynamics model from that transition. It then drew
simulated transitions from the learned model and ran value-iteration
backups on those simulations between every pair of real-environment
steps. The classical Dyna-Q instance kept a tabular maximum-likelihood
estimate of the transition and reward functions and used the simulated
backups to amortise the cost of real interaction. Dyna's lasting
contribution was conceptual rather than algorithmic. It made explicit
that the data the agent collects through interaction can be amplified
by a forecaster of the environment, that the forecaster's update and
the planner's update can run in the same outer loop, and that the
forecaster need not be a maximum-likelihood estimate even though the
original Dyna-Q chose it to be one. The classical caveat applied
already in 1990 and applies to every descendant. When the model is
inaccurate, the simulated transitions inject bias into the value
function, and the planner amplifies rather than corrects the error.
The value-aware programme of \S\ref{sec:vaml_ve} is the direct response
to that caveat at the loss-function level.

What the modern world-model literature adds is that the forecaster, the
value function, and the planner are trained jointly against the same
planning-relevant targets, which is the empirical face of the
value-equivalence principle of \S\ref{sec:vaml_ve}. Three architectural
families anchor the literature.

\citet{Schrittwieser2020MuZero} introduce the MuZero algorithm. A
representation network embeds the observation into a latent state $h_0$.
A dynamics network unrolls latent transitions
$h_{k+1}, r_{k+1} = g(h_k, a_k)$. A prediction network outputs a value
$v(h_k)$ and a policy prior $\pi(h_k)$ at each unrolled step. The three
networks are trained jointly against three planning-relevant targets at
every unrolled depth $k$, namely the empirical Bellman target for $v$,
the bootstrapped policy prior for $\pi$, and the realised reward for $r$.
No observation-reconstruction loss is used. The latent dynamics network
is therefore a value-equivalent model in the sense of
\citet{GrimmVALEQ2020} rather than a generative forecaster of pixels.
Embedded inside Monte Carlo Tree Search, the resulting agent matches
AlphaZero on Go, chess, and shogi without access to game rules, and
attains state of the art on $57$ Atari games without a simulator.
\citet{Antonoglou2022StochMuZero} extend the algorithm to stochastic
environments by replacing the deterministic latent dynamics with
afterstate-plus-chance-node representations, with matching performance on
$2048$ and backgammon.

A complementary line of work uses an explicit probabilistic model of the
dynamics together with model-predictive control. \citet{Chua2018PETS}
introduce probabilistic ensembles with trajectory sampling. The dynamics
model is an ensemble of probabilistic neural networks that captures both
aleatoric uncertainty, through the per-network output variance, and
epistemic uncertainty, through ensemble disagreement. At each control
step, candidate action sequences are scored by Monte Carlo rollouts
through the ensemble and the first action of the best sequence is
executed. The model-predictive-control step is a one-step rollout in the
sense of \citet{BertsekasRollout2022}, and the ensemble is the
probabilistic forecast that feeds it. The combination attains the
asymptotic performance of model-free baselines such as soft actor-critic
and proximal policy optimisation while using between $8$ and $125$ times
fewer environment samples on standard continuous-control tasks, which is
the cleanest empirical case for probabilistic-dynamics-plus-rollout on a
benchmark suite that an econometrician would recognise.

The third family is the Dreamer line of recurrent state-space models.
\citet{HafnerDreamer2020} train a recurrent latent dynamics model from
pixels, plan in the compact latent space rather than in observation
space, and update the actor by backpropagating analytic value gradients
through imagined trajectories. DreamerV3 \citep{HafnerDreamerV3_2025}
scales the recipe across more than $150$ control tasks under a single
hyperparameter configuration and is the first algorithm to collect
diamonds in Minecraft from pixels and sparse rewards without human
demonstrations.\footnote{The Dreamer architecture keeps an
observation-reconstruction term in the world-model loss, which makes it
not strictly value-equivalent. The reconstruction loss regularises the
latent representation in ways that the MuZero-style pure value-target
objective does not. The \citet{VoelckerCalib2025} calibration analysis of
\S\ref{sec:vaml_ve} applies most directly to the MuZero family and remains
an open question for Dreamer-style models.}

These architectures still treat the forecaster as a component that feeds
a separate planner or value function. The next section turns to the more
radical inversion, in which the forecaster is the policy itself.
